% Angepasste Seminarvorlage
%
\documentclass[12pt,a4paper,fancyheader]{scrartcl}

\usepackage[ngerman]{babel}			% deutsche Namen/Umlaute
\usepackage[utf8]{inputenc}			% Zeichensatzkodierung
\usepackage{fancyhdr}
\usepackage{graphicx}				% Einbinden von Bildern
\usepackage{color}					% Farben wenn es sein muß
\usepackage{hyperref}				% Klickbare Verweise und \autoref{label}
\usepackage{booktabs}				% "Schöne" Tabellen
\usepackage{amsmath}				% Mathematischer Formelsatz AMS
\usepackage{amsfonts}

%
%	Hier werden Titel, Bearbeiter und das Datum eingetragen
%
\newcommand\svthema{\LaTeX-Vorlage einer Semiararbeit}
\newcommand\svperson{Vorname Name, Vorname Name, Vorname Name}
\newcommand\svdatum{\today} % \today für das aktuelle Datum, sonst als Text eintragen
\newcommand\lvname{Modul: Modulname}
\newcommand\lvtyp{WS 2018}
\newcommand\lvinst{FOM Hochschule für Oekonomie \& Management}

% Das ist ein 100 Wörter umfassender Blindtext, um das Aussehen von längeren Texten zu testen.
\newcommand{\blindtext}{\textit{Lorem ipsum dolor sit amet, consetetur sadipscing elitr, sed diam nonumy eirmod tempor invidunt ut labore et dolore magna aliquyam erat, sed diam voluptua. At vero eos et accusam et justo duo dolores et ea rebum. Stet clita kasd gubergren, no sea takimata sanctus est Lorem ipsum dolor sit amet. Lorem ipsum dolor sit amet, consetetur sadipscing elitr, sed diam nonumy eirmod tempor invidunt ut labore et dolore magna aliquyam erat, sed diam voluptua. At vero eos et accusam et justo duo dolores et ea rebum. Stet clita kasd gubergren, no sea takimata sanctus est Lorem ipsum dolor sit amet.}}

\hypersetup{ % Thema und Author in die Meta-Daten der PDF
pdftitle={Seminararbeit zum Thema: \svthema} 
pdfauthor={\svperson}
}	
	
\begin{document}

% Hier wird der Titel-Bereich formatiert
% Die zuvor definierten Textbausteine werden hier verwendet.
\title{ \huge\textbf{\svthema} }
\author{ \textsl{\svperson} \\ \svdatum }
\date{ \normalsize \centering \includegraphics[width=0.3\textwidth]{FOM}\\\textsc{"`\lvname"'} \\ {\lvtyp} · {\lvinst} }
%eigene Veränderung- FOM Logo mit eingefügt auf dem Deckblatt

\maketitle
\thispagestyle{empty} % lässt die Seitennummer auf der Titelseite verschwinden
\pagenumbering{Roman}

\begin{abstract}
Hier steht eine kurze Zusammenfassung von circa 5--10 Zeilen.

\blindtext
\end{abstract}

\cleardoublepage
\tableofcontents			% Inhaltsverzeichnis ... nicht für kurze Dokumente!
\cleardoublepage
\listoffigures				% Abbildungsverzeichnis ... nicht für kurze Dokumente
\cleardoublepage
\pagenumbering{arabic}
\setcounter{page}{4}

\input{kapitel1.tex}

\input{kapitel2.tex}


% Das Literaturverzeichnis
\cleardoublepage
\bibliographystyle{IEEEtran}	% ieeetran verwenden, damit auch URLs angezeigt werden
\bibliography{seminar-lit}
\end{document}